\documentclass[12pt]{article}
\usepackage[utf8]{inputenc}
\usepackage[T1]{fontenc}
\usepackage{hyperref}
\title{Kaimb\'e Language Resources for K-12 Curriculum Development}
\author{Compiled by Claude Sonnet 5.5 (claude-sonnet-5-5) for Andr\'as Kornai}
\date{}
\begin{document}
\maketitle

\section{Introduction}
Kaimb\'e (ISO 639-3 xai, Glottocode kaim1235) is an extinct, unclassified language of Bahia, Brazil. It is very scarcely attested. About eight glosses were collected in 1961 by Wilbur Pickering from an elderly man in Massacar\'a (Euclides da Cunha, Bahia), who remembered a few words mixed with Kiriri. Several glosses are uncertain. Wikipedia gives an ethnic population of 1,135 (2020), which counts people, not speakers.

\section{Primary Educational Resources}
None found. No Kaimb\'e school materials or revitalization programme appeared in the searches.

\section{Linguistic and Reference Materials}
\begin{itemize}
\item Kaimb\'e language article, with the short word list (cites Meader 1978): \url{https://en.wikipedia.org/wiki/Kaimb%C3%A9_language}
\item Dissertation PDF from UFRJ (Pantheon), returned in the search; its relevance to Kaimb\'e was not checked: \url{https://pantheon.ufrj.br/bitstream/11422/23466/1/PTBSilva.pdf.pdf}
\end{itemize}

\section{Religious and Community Resources}
A Portuguese source says the Kaimb\'e originated from the Santíssima Trindade mission of Massacar\'a. No community organisation URL was seen.

\section{Digital and Audio Resources}
None found.

\section{Curriculum Development Context}
The documentation is too thin for a language-based curriculum. A heritage programme would have to rely on cultural and historical content, and on the revitalization experience of neighbouring Northeast Brazilian peoples such as the Pataxó (Patx\^oh\~a) and Kariri-Xoc\'o. The searches surfaced material on those, but it is not specific to Kaimb\'e: \url{https://periodicos.ufes.br/percursos/article/download/40923/27776}, \url{https://leetraindigena.ufscar.br/index.php/leetraindigena/article/download/43/33/153}.
\end{document}